\documentclass[11pt]{article}
\usepackage[a4paper,margin=2.3cm]{geometry}
\usepackage{amsmath,amssymb,amsthm}
\usepackage{mathtools}
\usepackage{empheq}
\usepackage{hyperref}
\usepackage{booktabs}
\usepackage{enumitem}
\usepackage{xcolor}
\usepackage{listings}
\usepackage{array}
\usepackage{tikz}
\usepackage{cleveref}

\hypersetup{colorlinks=true, linkcolor=blue!50!black, urlcolor=blue!60!black,
            citecolor=blue!50!black}

\newtheorem{proposition}{Proposition}
\newtheorem{lemma}{Lemma}
\newtheorem{corollary}{Corollary}
\newtheorem{theorem}{Theorem}
\theoremstyle{definition}
\newtheorem{definition}{Definition}
\newtheorem{remark}{Remark}

% --- Notation macros (kept consistent across Parts I and II) ---------------
\newcommand{\Tlock}{T-Lock}
\newcommand{\IRR}{\ensuremath{\mathrm{IRR}}}
\newcommand{\RF}{\ensuremath{\mathrm{RF}}}                 % RiskFactor
\newcommand{\FwdP}{\ensuremath{\mathrm{ForwardPrice}}}
\newcommand{\PVone}{\ensuremath{\mathrm{PV01}}}
\newcommand{\NDELTA}{\ensuremath{\mathrm{NDELTA}}}
\newcommand{\NPV}{\ensuremath{\mathrm{NPV}}}
\newcommand{\E}{\mathbb{E}}
\newcommand{\Q}{\mathbb{Q}}
\newcommand{\Drepo}{D^{\mathrm{repo}}}
\newcommand{\Dtl}{D^{\text{T-Lock}}}
\newcommand{\Tdel}{T_{\mathrm{del}}}
\newcommand{\Adel}{A_{\mathrm{del}}}
\newcommand{\Ptl}{P^{\text{T-Lock}}}
\newcommand{\ytl}{y^{\text{T-Lock}}}
\newcommand{\Bf}{B_f}
\newcommand{\rrepo}{r^{\mathrm{repo}}}
\newcommand{\rfun}{r^{\mathrm{fun}}}
% Derivative-with-respect-to macros (k-th order)
\newcommand{\dDy}[1]{D_{y}^{(#1)}}              % k-th y-derivative
\newcommand{\dDP}[1]{D_{P}^{(#1)}}              % k-th P-derivative
\DeclareMathOperator{\OCI}{OCI}

\title{\textbf{The Treasury Lock}\\[4pt]
       \large A comprehensive study\\[2pt]
       \normalsize Part I -- Practitioner Pricing Model \quad|\quad
                  Part II -- Hedging and Convexity Analysis (Pucci 2019)}
\author{Working notes}
\date{}

\begin{document}
\maketitle
\vspace{-2.0em}

\begin{center}
\small\itshape
Source papers:\\
(1) \emph{Treasury Lock Model}, anonymous practitioner note (no author, no date).\\
(2) M.\ Pucci, \emph{Hedging the Treasury Lock}, Banca IMI SpA, Milan,
August 19, 2019, SSRN no.\ 3386521,
\url{https://ssrn.com/abstract=3386521}.
\end{center}

\hrule
\bigskip

\tableofcontents
\bigskip
\hrule

\section{Introduction and Trade Description}\label{sec:intro}

A \emph{Treasury Lock} (\Tlock) is a customised over-the-counter agreement
that fixes either the \emph{yield} or the \emph{price} of a specified
Treasury security for a specific period. The buyer is protected from a rise
in the yield (equivalently a fall in price) of the underlying Treasury
during the lock period. The lock period is normally one week to twelve
months, most often $\le 6$ months.

\Tlock s are most commonly used as a \emph{pre-hedge} for an upcoming USD
bond or loan issuance: the driver of the fixed rate paid by the borrower
is the sum of a credit spread and the benchmark Treasury yield, so the
client locks the Treasury component a few months ahead of issue.

Being \textbf{long} a \Tlock\ at strike $L$ (the Locked Yield) means
receiving at Expiry $\Tdel$ an amount
\begin{equation}
g \;=\; a \cdot N \cdot \RF_{\Tdel} \cdot \bigl( \IRR_{\Tdel} - L \bigr),
\qquad a = +1 \;(\text{long}),\; -1 \;(\text{short}),
\label{eq:payoff-pucci}
\end{equation}
where
\begin{itemize}[leftmargin=1.4em,itemsep=2pt,topsep=2pt]
\item $N$ is the notional;
\item $\IRR_{\Tdel}$ is the bond's internal rate of return at expiry;
\item $\RF_t = -\partial P_t(y)/\partial y\big|_{y=\IRR_t}$ is minus the
      derivative of price w.r.t.\ yield (so $\RF = -10{,}000\cdot\text{DV01}$).
\end{itemize}

Two contract variants exist:
\begin{itemize}[leftmargin=1.4em,itemsep=2pt,topsep=2pt]
\item \textbf{Standard \Tlock}: the underlying Treasury is specified at trade
      inception.
\item \textbf{Then-Current \Tlock}: the underlying is the on-the-run $n$-Year
      benchmark prevailing at $\Tdel$ ($n = 1,\dots,30$; typically $n=5,7,10$).
\end{itemize}

\textbf{Hedge accounting.} The \Tlock\ can be designated as the hedging
instrument in a cash-flow hedge of an anticipated borrowing transaction.
P\&L is set aside in $\OCI$ until expiry; if the borrowing materialises
the P\&L is released amortised over the bond's life, partially offsetting
debt-service cash flows. If financial close fails, the full P\&L recycles
as cost/income at expiry (IFRS~9).

This document brings together two complementary perspectives:

\begin{description}[leftmargin=1.6em,itemsep=2pt]
\item[Part I --- Practitioner Pricing Model (anonymous note).]
Gives the operational pricing formulas a trading desk needs: the bond
forward price with explicit coupon stripping (\cref{eq:bond-forward}),
and three \NPV\ formulas for the three lock conventions
(clean-price lock, yield lock, dirty-price lock).
\item[Part II --- Hedging and Convexity Analysis (Pucci 2019).]
Justifies the desk practice of \emph{booking} a (short) \Tlock\ as a (long)
forward contract on the same security. Shows the resulting overhedge
is (i)~bounded analytically by a Taylor remainder, (ii)~generates
non-negative gamma P\&L (negative-convexity overhedge), (iii)~statistically
negligible at typical $\le 6$M tenors with at-the-money strikes.
\end{description}


\section{Unified Notation}\label{sec:notation}

\begin{center}
\renewcommand{\arraystretch}{1.12}
\begin{tabular}{l p{12cm}}
\toprule
\textbf{Symbol} & \textbf{Meaning} \\
\midrule
$t$, $T_0$ & Valuation date / spot settlement date \\
$\Tdel$, $t_e$ & Delivery date $=$ \Tlock\ Expiry (interchangeable) \\
$P_0$ & Spot clean price of the underlying bond \\
$A_0$, $\Adel$ & Accrued interest at $T_0$ and at $\Tdel$ \\
$C_i,\ T_i$ & $i$-th coupon and its payment date, $T_0 < T_i \le \Tdel$ \\
$c$ & Coupon rate of the bond (continuous-compounding convention, Part II) \\
$\alpha_i$ & Accrual factor for the $i$-th coupon period \\
$N$ & Number of coupons paid in $(T_0, \Tdel]$ (Part I)\ /\ bond notional (Part II) \\
$M$ & Trade amount for a yield-locked \Tlock\ (Part I, \cref{eq:npv-yield}) \\
$P^{\mathrm{mkt}}_t,\ P_t(y)$ & Market dirty price /\ price as a function of yield \\
$\IRR_t$ & Internal rate of return at $t$: solution of $P^{\mathrm{mkt}}_t = P_t(\IRR_t)$ \\
$\Bf$ & Bond forward (clean) price at delivery \\
$y_f$ & Forward bond yield (= the $\IRR$ implied by $\Bf$) \\
$L,\ \Ptl,\ \ytl$ & Locked yield / locked clean price / locked yield (strike of the \Tlock) \\
$\Drepo(s, T)$ & Discount factor from $T$ back to $s$, repo curve \\
$\Dtl(s, T)$ & Discount factor from $T$ back to $s$, bond zero-yield curve \\
$D_{tT}$ & Risk-free zero-coupon bond from $t$ to $T$ (Part II) \\
$\rrepo,\ \rfun$ & Repo rate /\ unsecured funding rate \\
$\RF_t$ & Risk factor: $-\partial P_t/\partial y$ at $y = \IRR_t$ \\
$\PVone(y_f)$ & Centred-difference forward-yield risk factor,
                $B(y_f + \tfrac12\text{bp}) - B(y_f - \tfrac12\text{bp})$ \\
$a$ & Sign convention: $a=+1$ long, $a=-1$ short \\
\bottomrule
\end{tabular}
\end{center}

\textbf{A note on conventions.} The anonymous practitioner paper (Part I)
uses simple compounding inside the bond forward formula and a centred
finite-difference $\PVone$. Pucci (Part II) uses the
continuous-compounding bond price
\begin{equation}
P_t(y) \;=\; e^{-(t_n - t)y} + c\sum_{t_i > t} \alpha_i\, e^{-(t_i - t)y}
\label{eq:bond-cont}
\end{equation}
to derive analytical greeks. We keep the two conventions in their
respective Parts (this matches the source) and reconcile them
in \cref{sec:validation}.

\bigskip\hrule\medskip
\begin{center}
\Large\bfseries Part I --- Practitioner Pricing Model
\end{center}
\hrule\bigskip

\section{The Bond Forward Pricing Building Block}\label{sec:bondfwd}

The pricing of every \Tlock\ NPV formula in this Part reduces to the
\emph{forward clean price} of the underlying bond at delivery.
Ignoring bid-offer (the model is mid-market), the bond forward clean
price is\footnote{This is the standard cash-and-carry no-arbitrage
identity: buy the bond spot for $P_0+A_0$, finance at the repo rate,
collect intermediate coupons, deliver at $\Tdel$.}
\begin{empheq}[box=\fbox]{equation}
\Bf \;=\;
\frac{\,P_0 + A_0\,}{\Drepo(T_0,\Tdel)}
\;-\;
\sum_{i=1}^{N}\frac{C_i}{\Drepo(T_i,\Tdel)}
\;-\;\Adel,
\label{eq:bond-forward}
\end{empheq}
where
\begin{itemize}[leftmargin=1.4em,itemsep=1pt]
\item $P_0$ = clean price of the underlying bond at spot;
\item $A_0,\,\Adel$ = accrued interest on a unit of the bond at $T_0$
      and at $\Tdel$;
\item $C_i$ = the $i$-th coupon paid at $T_i$, $T_0 < T_i \le \Tdel$;
\item $\Drepo(s,T)$ = discount factor from $T$ back to $s$, retrieved
      from the repo zero curve;
\item $N$ = number of coupons paid in $(T_0, \Tdel]$.
\end{itemize}

\textbf{Pricing parameter \texttt{BOND\_FROM\_QUOTE}.} By default this
flag is \texttt{true}: $P_0 = $ \texttt{MarketQuote} if available; else
$P_0$ is derived from yield via the price function $P_0(y)$. This affects
sensitivity computation (whether $P_0$ moves with bond shifts).

The forward yield $y_f$ is implicitly defined by
\begin{equation}
P_{\Tdel}\bigl(y_f\bigr) \;=\; \Bf,
\label{eq:forward-yield}
\end{equation}
where $P_{\Tdel}(y)$ is the bond's price-yield function evaluated
\emph{as if it settles on $\Tdel$}.


\section{The Three Lock Conventions}\label{sec:three-conv}

A \Tlock\ can be struck on any of three quantities. Throughout this
section the present value is for the \emph{investor / buyer / long}.

\subsection{Lock the clean price ($\Ptl$ struck)}

\begin{empheq}[box=\fbox]{equation}
\NPV^{\Tlock} \;=\; N \cdot \bigl( \Ptl - \Bf \bigr) \cdot
\Dtl(t,\Tdel)
\label{eq:npv-price}
\end{empheq}
where
\begin{itemize}[leftmargin=1.4em,itemsep=1pt]
\item $N$ = quantity of the underlying bond;
\item $\Ptl$ = locked clean price of the underlying bond (the strike);
\item $\Bf$ = bond forward clean price, \cref{eq:bond-forward};
\item $\Dtl(t,\Tdel)$ = discount factor from $\Tdel$ back to valuation
      date $t$, retrieved from the \emph{bond zero-yield curve}.
\end{itemize}

\noindent The two-curve structure is essential: $\Bf$ is built on the
repo curve (carry of the security), while the final payoff is discounted
on the bond zero-yield curve (the relevant funding for the cash payoff).

\subsection{Lock the yield ($\ytl$ struck)}

\begin{empheq}[box=\fbox]{equation}
\NPV^{\Tlock} \;=\; M \cdot \bigl( y_f - \ytl \bigr) \cdot
\PVone(y_f) \cdot \Dtl(t,\Tdel)
\label{eq:npv-yield}
\end{empheq}
where
\begin{itemize}[leftmargin=1.4em,itemsep=1pt]
\item $M$ = trade amount for a yield lock;
\item $y_f$ = forward bond yield, \cref{eq:forward-yield};
\item $\ytl$ = locked bond yield (the strike);
\item $\PVone(y_f)$ is the \emph{centred-difference forward-yield risk
      factor}:
\begin{equation}
\PVone(y_f) \;=\; B\bigl(y_f + 0.5\text{bp}\bigr)
                 - B\bigl(y_f - 0.5\text{bp}\bigr),
\label{eq:pv01-fwd}
\end{equation}
i.e.\ the price change of a unit bond per 1bp shift of $y_f$;
\item $\Dtl(t,\Tdel)$ as above.
\end{itemize}

\begin{remark}\label{rem:pv01-vs-fwd}
$\PVone(y_f)$ \emph{is not} the conventional spot $\PVone$. It is the
risk-factor of the \emph{forward-settled} bond, evaluated at the
forward yield $y_f$. Confusing the two is a frequent implementation
bug.
\end{remark}

\subsection{Lock the dirty price}

The dirty price is the sum of the clean price and accrued interest for
a unit bond. The model first computes the implied clean price strike
\[
\Ptl \;=\; \text{LockedDirtyPrice} - \Adel
\]
and then applies \cref{eq:npv-price}. No new formula is required.


\section{Risks and Sensitivities (Practitioner View)}\label{sec:risks-practitioner}

The anonymous note explicitly cautions:
\begin{quote}\itshape
There is no specific indicator of risk sensitivity for \Tlock s. Both
$\PVone$ and $\NDELTA$ are underlier-related. We apply them here as
approximations of the \Tlock\ sensitivity. Therefore we consider it a
limitation to apply them as a sensitivity indicator for \Tlock s.
\end{quote}

The natural risk sensitivities the desk reports are therefore:
\begin{itemize}[leftmargin=1.4em,itemsep=2pt]
\item $\PVone$ of the \emph{underlying bond}, multiplied by quantity
      and discount factor;
\item $\NDELTA$ ($n$th-order delta) of the underlying bond, similarly
      scaled;
\item Repo-bucket sensitivity, since \cref{eq:bond-forward} depends on
      $\Drepo$ at multiple coupon nodes.
\end{itemize}
Part II will sharpen these into formal greeks of the \Tlock\ itself
and explain why the practitioner proxy is an \emph{overhedge}, not a
sensitivity match.


\bigskip\hrule\medskip
\begin{center}
\Large\bfseries Part II --- Hedging and Convexity Analysis
\end{center}
\hrule\bigskip

\section{The Arbitrage-Free Pricing Formula}\label{sec:af-price}

Recall the spot price-yield identity:
\begin{equation}
P^{\mathrm{mkt}} \;=\; P(\IRR),
\label{eq:irr-def}
\end{equation}
i.e., the $\IRR$ is the yield that reprices the bond.

The \Tlock\ prescribes the compounded-interest formula
\begin{equation}
P_t(y) \;=\; \prod_{i=1}^n \frac{1}{1+\alpha_i y}
\;+\; c \sum_{i=1}^n \alpha_i \prod_{j=1}^i \frac{1}{1+\alpha_j y},
\label{eq:bond-comp}
\end{equation}
with coupon rate $c$, issue date $t_0 \ge t$, coupon dates
$t_1<\dots<t_n$ and accrual factors $\alpha_1,\dots,\alpha_n$.
When $t$ lies in a coupon period $(t_0,t_1]$ the formula generalises by
multiplying with the stub discount $\left(\frac{1}{1+\alpha_1 y}\right)^{(t_1-t)/(t_1-t_0)}$.

With $\RF_t = -\partial P_t/\partial y\big|_{y=\IRR_t}$, the \Tlock\
payoff at $t_e$ is $g$ in \cref{eq:payoff-pucci}.

\textbf{First-order interpretation.} At first order in
$(\IRR_{t_e} - L)$:
\begin{align}
g &= -a\cdot \frac{\partial P_{t_e}(y)}{\partial y}\bigg|_{y=\IRR_{t_e}}
       \cdot (\IRR_{t_e} - L) \notag\\
  &\sim a\cdot \bigl( P_{t_e}(L) - P_{t_e}(\IRR_{t_e}) \bigr)
   \;=\; a\cdot \bigl( P_{t_e}(L) - P^{\mathrm{mkt}}_{t_e} \bigr).
\label{eq:first-order}
\end{align}
The \Tlock\ thus measures the gap between the market price of the
bond at expiry and its price with cash-flows discounted at the locked
rate $L$. We will show this identifies an \emph{over}hedge for
$a=+1$ and an \emph{under}hedge for $a=-1$.

\textbf{General formula.} In an arbitrage-free framework
\begin{empheq}[box=\fbox]{equation}
v_t \;=\; a \cdot D_{t t_e} \cdot
        \E^{t_e}_t\!\bigl[\, \RF_{t_e}\cdot(\IRR_{t_e} - L)\,\bigr],
\label{eq:af-pricing}
\end{empheq}
where the expectation is under the $t_e$-forward measure and $D_{t t_e}$
is the risk-free $t_e$-zero-coupon bond seen at $t$. Default risk is
ignored: short tenor + US-government issuer.

\textbf{Why this is not just a CMT swaplet.} If $\RF_{t_e}$ is
\emph{frozen}, the payoff superficially resembles a CMS-let on the
bond curve (a Constant Maturity Treasury swaplet, cf.\ Pucci 2014,
SSRN 3387961). This interpretation is \emph{misleading} because
freezing $\RF$ overshadows --- and actually reverses --- a key
feature: while a CMS-let has positive convexity, the \Tlock\ has
\emph{negative} convexity locally near $L$
(see \cref{sec:gamma}).


\section{The Forward-Contract Proxy and Overhedge}\label{sec:overhedge}

\subsection{The proxy strike}

For ease of analytics, from now on we use the continuous-compounding
bond price \cref{eq:bond-cont}:
\[
P_t(y) \;=\; e^{-(t_n-t)y} + c\sum_{t_i>t}\alpha_i\, e^{-(t_i-t)y}.
\]

From \cref{eq:first-order}, the payoff of a long Standard \Tlock\ struck
at $L$ is approximated, at first order, by the $t$-value of a
\emph{short} forward contract struck at $P_{t_e}(L)$. The proxy strike is
\begin{align}
K &= P_{t_e}(L) \label{eq:K-exact}\\
  &\sim \FwdP_t(t_e) + \RF_{t_e}\cdot(\IRR_t - L), \label{eq:K-approx}
\end{align}
where $\IRR_t$ is the $t_e$-forward $\IRR$ defined implicitly by
$P_{t_e}(\IRR_t) = \FwdP_t(t_e)$.

The forward proxy is a convenient booking representation when the front
office system does not feature a \Tlock\ typology.

\subsection{The overhedge error is the Taylor remainder}

Substituting $P$ by its first-order Taylor polynomial about $L$,
\cref{eq:first-order} approximation becomes an exact overhedge once
we identify it with the Taylor remainder.

\begin{proposition}[Overhedge bound]\label{prop:overhedge}
Let
\[
R_1(y) \;=\; P(y) - P(L) - D_y[P](L)\cdot(y-L).
\]
Then $R_1$ has a unique stationary point at $y=L$ which is also a
maximum (since the second derivative of $R_1$ is negative everywhere).
Moreover, by the Remainder Estimation Theorem,
\begin{equation}
|R_1(y)| \;\le\; \frac{M\,|y - L|^2}{2},
\label{eq:rem-bound}
\end{equation}
where $M = \max_{y\in I} |\dDy{2}[P](y)| = \dDy{2}[P](\min I)$ (the maximum
is attained at the left-most point of the interval $I$ because
$\dDy{3}[P] < 0$).
\end{proposition}

\begin{proof}
$\dDy{2}[P]>0$ (positive convexity, \cref{eq:Dyy} below), so $R_1$ is
concave on $I$. Its first derivative
$D_y[R_1](y) = D_y[P](y) - D_y[P](L)$
vanishes only at $y=L$, giving the unique critical point.
Concavity makes it a maximum. The remainder bound is the standard
Taylor result with $|D_y^{(2)}[P]|$ maximised on $I$, attained at the
left endpoint because $D_y^{(3)}[P]<0$.
\end{proof}

\textbf{Numerical magnitude.} For the 10Y benchmark
\[
M \;\sim\; (t_n - t)^2 + c\sum_{t_i>t}(t_i-t)^2\,\alpha_i \;\sim\; 150
\]
(taking the extreme case $y=0$; minimum 1Y CMT yield was 0.08\% in
September 2011, so this is a hard upper bound). A 30\% yield move ---
an upper bound for a 6M tenor --- yields a dollar error of
$\approx 0.005\%$ of notional. The overhedge is therefore tight.


\section{Greeks of the \Tlock}\label{sec:greeks}

\subsection{Derivatives of the bond price}\label{sec:bond-derivs}

From the continuous-compounding form \cref{eq:bond-cont}:
\begin{align}
D_y[P_t]   &= -(t_n-t)\,e^{-(t_n-t)y}
              - c\sum_{t_i>t}\alpha_i(t_i-t)\,e^{-(t_i-t)y}
              \;<\;0,
              \label{eq:Dy}\\[2pt]
\dDy{2}[P_t]&= (t_n-t)^2 e^{-(t_n-t)y}
              + c\sum_{t_i>t}\alpha_i(t_i-t)^2 e^{-(t_i-t)y}
              \;>\;0,
              \label{eq:Dyy}\\[2pt]
\dDy{3}[P_t]&= -(t_n-t)^3 e^{-(t_n-t)y}
              - c\sum_{t_i>t}\alpha_i(t_i-t)^3 e^{-(t_i-t)y}
              \;<\;0.
              \label{eq:Dyyy}
\end{align}
The convexity of the bond is $\dDy{2}[P]/P > 0$.

\subsection{Payoff derivatives}

Writing the payoff as a function of yield, $g(y) = -a\,D_y[P](y-L)$,
\begin{align}
D_y[g] &= -a\bigl[\dDy{2}[P]\cdot(y-L) + D_y[P]\bigr],
         \label{eq:Dy-g}\\
\dDy{2}[g] &= -a\bigl[\dDy{3}[P]\cdot(y-L) + 2\dDy{2}[P]\bigr].
         \label{eq:Dyy-g}
\end{align}

\subsection{Delta of the \Tlock}

For $a=+1$ and $|y-L|$ small enough (which is the typical business case:
ATM strike, short tenor),
\begin{equation}
D_y[g] > 0.
\label{eq:delta-y}
\end{equation}
The delta with respect to the security price is then
\begin{equation}
D_P[g] \;=\; D_y[g]\cdot D_P[y] \;<\; 0,
\label{eq:delta-P}
\end{equation}
consistent with a long \Tlock\ being equivalent to shorting the
security.

A sharp characterisation: solving \cref{eq:Dy-g} for $L$,
\begin{equation}
D_y[g]>0\ \text{and}\ \Delta<0
\iff
L > y + \frac{D_y[P]}{\dDy{2}[P]}.
\label{eq:delta-sign}
\end{equation}
Equivalently $y - L < -D_y[P]/\dDy{2}[P] > 0$: as long as $y$ does not
overshoot $L$, $\Delta$ stays negative.

\textbf{Hedging.} Accounting for discounting, a trader who books the
\Tlock\ short via the proxy forward of \cref{sec:overhedge} can
delta-hedge by \emph{shorting} an amount of bond equal to
\begin{equation}
\Delta_{\text{hedge}} \;\sim\; D_{t t_e}\cdot D_P[g].
\label{eq:delta-hedge}
\end{equation}
The hedge is continuously rebalanced.

\subsection{Gamma of the \Tlock}\label{sec:gamma}

From the second-derivative chain rule (see \cref{sec:simple-math},
\cref{eq:chain2}):
\begin{align}
\dDP{2}[g]
  &= \dDy{2}[g]\cdot(D_P[y])^2 + D_y[g]\cdot\dDP{2}[y]
  \notag \\
  &= \dDy{2}[g]\cdot\bigl(1/D_y[P]\bigr)^2
   - D_y[g]\cdot\bigl[\dDy{2}[P]/(D_y[P])^3\bigr].
\label{eq:gamma-chain}
\end{align}
Substituting \cref{eq:Dy-g,eq:Dyy-g} into \cref{eq:gamma-chain}, for
$|y-L|$ small:
\begin{align}
\dDP{2}[g]
  &\sim -a\bigl[2\,\dDy{2}[P]\bigr]\bigl(1/D_y[P]\bigr)^2
        + a\bigl[D_y[P]\bigr]\bigl[\dDy{2}[P]/(D_y[P])^3\bigr] \notag\\
  &=    -a\bigl[2\,\dDy{2}[P]\bigr]\bigl(1/D_y[P]\bigr)^2
        + a\,\bigl[\dDy{2}[P]/(D_y[P])^2\bigr] \notag\\
  &=    -a\cdot\dDy{2}[P]\cdot\bigl(1/D_y[P]\bigr)^2.
\label{eq:gamma-local}
\end{align}

For $a=+1$ and $|y-L|$ small, $\dDy{2}[P]>0$ and $(D_y[P])^{-2}>0$, so
\[
\boxed{\dDP{2}[g] < 0 \quad\text{(negative gamma / negative convexity).}}
\]
\textbf{Practical consequence.} Dynamically replicating a long \Tlock\
under no-volatility assumption is an \emph{over}hedge generating
non-negative P\&L for the trader (cf.\ Carr-Madan 2002 [14, Eq.~9]).
Equivalently: delta-hedging a \emph{short} \Tlock\ position generates
\emph{gamma profits}.

\subsection{Sharp negativity bounds}

A non-local characterisation: from \cref{eq:gamma-chain},
\begin{equation}
\Gamma < 0
\iff
(y-L)\!\cdot\!\Bigl[(\dDy{2}[P])^2 - D_y[P]\cdot\dDy{3}[P]\Bigr]
> D_y[P]\cdot\dDy{2}[P].
\label{eq:gamma-bound-raw}
\end{equation}
The bracket on the LHS is strictly negative:
\begin{lemma}\label{lem:bracket}
$(\dDy{2}[P])^2 - D_y[P]\cdot\dDy{3}[P] < 0$.
\end{lemma}
The proof is deferred to \cref{sec:more-math} (\cref{eq:bracket-proof}).
Given the lemma, dividing \cref{eq:gamma-bound-raw} by the (negative)
bracket flips the inequality:
\begin{equation}
\Gamma < 0
\iff
y - L \;<\;
\frac{D_y[P]\cdot\dDy{2}[P]}{(\dDy{2}[P])^2 - D_y[P]\cdot\dDy{3}[P]}
\;>\; 0.
\label{eq:gamma-bound}
\end{equation}
Again: as long as $y$ does not overshoot $L$, $\Gamma$ stays negative.
\Cref{eq:gamma-bound} gives an explicit \textbf{negativity-of-gamma
upper bound} we will estimate empirically in \cref{sec:stats}.


\section{Replication via Repo}\label{sec:repo}

The bond underlying the proxy forward is most often shorted via a
\emph{reverse repo} trade. \Cref{eq:delta-P} gives the delta with respect
to the forward price, which is interpreted as the size of the repo to
trade.

The forward price relation (zero-haircut) is
\begin{empheq}[box=\fbox]{equation}
\FwdP_t(t_e) \;=\;
P^{\mathrm{mkt}}_t \bigl(1 + \rrepo_{t t_e}\cdot(t_e - t)\bigr)
\;-\; c\sum_{t<t_i \le t_e}\alpha_i
        \bigl(1 + \rrepo_{t_i t_e}\cdot(t_e - t_i)\bigr).
\label{eq:fwd-repo}
\end{empheq}
This is the simple-compounding equivalent of \cref{eq:bond-forward}.
We assume zero haircut throughout: a good approximation for Treasuries.

\textbf{Repo-rate sensitivity.} By the chain rule on \cref{eq:fwd-repo}
and \cref{eq:K-approx},
\begin{equation}
\frac{\partial g(\rrepo)}{\partial \rrepo}
\;=\;
\frac{\partial g(y)}{\partial y}\Big|_{y=\IRR_t}
\cdot
\frac{\partial y(P)}{\partial P}\Big|_{P=\FwdP_t}
\cdot
\frac{\partial \FwdP(\rrepo)}{\partial \rrepo}.
\label{eq:repo-sensitivity}
\end{equation}
A higher repo rate $\Rightarrow$ higher forward price $\Rightarrow$ lower
forward yield $\Rightarrow$ \emph{cheaper} long \Tlock\ position.

\textbf{Specialness.} For benchmark Treasuries trading "on special",
the actual repo rate may be significantly below GC. Historical extreme:
during a liquidity squeeze the 10Y repo dropped to $-300$bp
(Leong 2014). The fair value of the long \Tlock\ rises accordingly.

\textbf{EUR aside.} Negative EUR repo rates make a hypothetical
\Tlock\ on a European benchmark more expensive than a USD one.

\begin{center}\small
\begin{tabular}{lcc}
\toprule
\textbf{Maturity} & \textbf{Bid} & \textbf{Ask} \\
\midrule
1M & 2.57\% & 2.52\% \\
3M & 2.60\% & 2.55\% \\
6M & 2.63\% & 2.58\% \\
12M & 2.70\% & 2.65\% \\
\bottomrule
\end{tabular}
\\[1pt]
\textit{Table: GC repo rates for the 10Y as of 24 Jan 2019.}
\end{center}


\section{The Then-Current \Tlock\ and Roll P\&L}\label{sec:roll}

In a Then-Current \Tlock\ the underlying is determined at $t_e$ as the
on-the-run benchmark prevailing at $t_e$. At inception this security
is unknown / unlisted in the position-keeping system.

\textbf{Booking convention.} The trader books the trade as a
\emph{Standard \Tlock\ with the currently on-the-run security}.
At each benchmark roll date $t_{\mathrm{roll}}$, the booking is
revised: the underlying is switched to the new on-the-run, generating
a \emph{roll P\&L}.

\subsection{The roll P\&L formula}

To keep notation light, assume the security is switched from a bond
with coupon $c$ to a bond with coupon $\hat c$, \emph{all other details
equal}. The NPV jumps from
\begin{align}
\pi_t  &= D_{t t_e}\bigl[K - F_t(t_e)\bigr], \label{eq:roll-pi}\\
\hat\pi_t &= D_{t t_e}\bigl[\hat K - \hat F_t(t_e)\bigr], \label{eq:roll-pihat}
\end{align}
with $K = P_{t_e}(L)$, $\hat K = \hat P_{t_e}(L)$, and the new
price-yield curve
\begin{equation}
\hat P_t(y) \;=\; e^{-(t_n-t)y}
                 + \hat c\sum_{t_i>t}\alpha_i e^{-(t_i-t)y}.
\label{eq:Phat}
\end{equation}
The P\&L due to the switch is
\begin{equation}
\hat\pi_t - \pi_t
\;=\; D_{t t_e}\Bigl[(\hat K - K) - (\hat F_t(t_e) - F_t(t_e))\Bigr].
\label{eq:roll-pl}
\end{equation}
The strike gap is purely coupon-driven:
\begin{equation}
\hat K - K \;=\; (\hat c - c)\sum_{t_i > t_e}\alpha_i\,e^{-(t_i-t_e)L}.
\label{eq:strike-gap}
\end{equation}
By definition $F_t(t_e) = P_{t_e}(R_t)$ and
$\hat F_t(t_e) = \hat P_{t_e}(\hat R_t)$, so
\begin{align}
\hat F_t - F_t
  &= \hat P_{t_e}(\hat R_t) - P_{t_e}(R_t) \notag\\
  &= \bigl[\hat P_{t_e}(\hat R_t) - P_{t_e}(\hat R_t)\bigr]
    + \bigl[P_{t_e}(\hat R_t) - P_{t_e}(R_t)\bigr] \notag\\
  &= (\hat c - c)\!\sum_{t_i > t_e}\!\alpha_i e^{-(t_i-t_e)\hat R_t}
    + \bigl[P_{t_e}(\hat R_t) - P_{t_e}(R_t)\bigr].
\label{eq:fwd-gap}
\end{align}

\textbf{Short-tenor approximation.} When $t_n - t_e$ is small (the bond
at $t_e$ is close to maturity), the last bracket is linearisable:
\begin{equation}
P_{t_e}(\hat R_t) - P_{t_e}(R_t)
\;\sim\;
\Bigl[(t_n-t_e) + c\!\sum_{t_i>t_e}\!\alpha_i(t_i-t_e)\Bigr]
\bigl(R_t - \hat R_t\bigr).
\label{eq:short-tenor}
\end{equation}
Combining \cref{eq:strike-gap,eq:fwd-gap,eq:short-tenor} into
\cref{eq:roll-pl}, the (forward) P\&L satisfies
\begin{empheq}[box=\fbox]{align}
\frac{\hat\pi_t - \pi_t}{D_{t t_e}}
&\sim
(\hat c - c)\bigl(\hat R_t - L\bigr)
\sum_{t_i>t_e}\alpha_i(t_i-t_e) \notag\\
&\quad + (\hat R_t - R_t)\Bigl[(t_n-t_e) + c\!\sum_{t_i>t_e}\!\alpha_i(t_i-t_e)\Bigr].
\label{eq:roll-PL}
\end{empheq}
A symmetric expansion gives
\begin{equation}
\frac{\hat\pi_t - \pi_t}{D_{t t_e}}
\sim (c-\hat c)(L - R_t)\sum_{t_i>t_e}\alpha_i(t_i-t_e)
   + (\hat R_t - R_t)\Bigl[(t_n-t_e) + \hat c\sum_{t_i>t_e}\alpha_i(t_i-t_e)\Bigr].
\label{eq:roll-PL-sym}
\end{equation}
Hence the security change generates little forward-contract P\&L if
\begin{enumerate}[leftmargin=1.6em,itemsep=1pt]
\item $t_n - t_e$ is small (close to bond maturity),
\item $L$ is close to $R_t$,
\item there is no significant $\IRR$ gap between the two securities.
\end{enumerate}

\textbf{General (longer-tenor) version.} Setting
\[
A_t(y) \;\equiv\; \sum_{t_i > t} \alpha_i\, e^{-(t_i-t)y},
\]
removing the "close-to-maturity" restriction yields
\begin{empheq}[box=\fbox]{equation}
\frac{\hat\pi_t - \pi_t}{D_{t t_e}}
\;\sim\;
(\hat c - c)\cdot D_y[A_{t_e}]\big|_{y=\hat R_t}\cdot(L - \hat R_t)
\;+\; D_y[P_{t_e}]\big|_{y=\hat R_t}\cdot(R_t - \hat R_t),
\label{eq:roll-general}
\end{empheq}
with
\begin{equation}
D_y[A_t] \;=\; -\sum_{t_i > t}\alpha_i(t_i-t)\,e^{-(t_i-t)y}.
\label{eq:DyA}
\end{equation}

\subsection{Forming the new benchmark}\label{sec:auction}

US Treasury benchmarks are reissued periodically by auction. The
auction process for a roll proceeds:
\begin{enumerate}[leftmargin=1.6em,itemsep=1pt]
\item Evaluate $R_t$, the IRR of the current benchmark.
\item Locate the prospected coupon $\hat c$ so the new issue is at par:
      $P^{(\hat c)}_t(R_t) = 1$.
\item The auction filters market appetite, producing a final coupon
      $\hat c'$ (possibly $\ne \hat c$).
\item Imply the new $\IRR$ from $P^{(\hat c')}_t(\hat R_t) = 1$.
\end{enumerate}
If the market confirms "at the old level" we land at $\hat R_t = R_t$,
which by \cref{eq:roll-general} minimises the roll P\&L.

\subsection{Strike gap for the Then-Current proxy}

If at expiry the new benchmark coincides with the issue date, the
proxy forward strike at inception is
\begin{equation}
K \;=\; 1 + \RF_{t_e}\cdot(R_t - L),
\label{eq:K-then-current}
\end{equation}
where $\RF_{t_e}$ is computed on the \emph{proxy} security.
At $t_{\mathrm{roll}}$ the strike gap becomes
\begin{align}
\text{StrikeGap}
&= \RF^{\mathrm{new}}_{t_e}\cdot(\hat R_{t_{\mathrm{roll}}} - L)
 - \RF^{\mathrm{old}}_{t_e}\cdot(R_{t_{\mathrm{roll}}} - L) \notag\\
&= \bigl(\RF^{\mathrm{new}} - \RF^{\mathrm{old}}\bigr)\cdot
   \bigl(\hat R_{t_{\mathrm{roll}}} - L\bigr)
   - \RF^{\mathrm{old}}\cdot
   \bigl(R_{t_{\mathrm{roll}}} - \hat R_{t_{\mathrm{roll}}}\bigr).
\label{eq:strike-gap-tc}
\end{align}


\section{Statistical Validation (Pucci 2019)}\label{sec:stats}

These tables are reproduced verbatim from Pucci 2019. They define
the empirical envelopes against which a pricing library can be
benchmarked.

\subsection{IRR descriptive statistics (since Jan-2008)}
\begin{center}\small
\begin{tabular}{lccc}
\toprule
\textbf{IRR} & \textbf{2Y} & \textbf{5Y} & \textbf{10Y}\\
\midrule
Average & 0.88\% & 1.61\% & 2.40\%\\
StDev & 0.75\% & 0.61\% & 0.52\%\\
5th pct & 0.24\% & 0.70\% & 1.62\%\\
95th pct & 2.59\% & 2.77\% & 3.34\%\\
\bottomrule
\end{tabular}
\end{center}

\subsection{Overhedge error (6M T-Lock on 10Y)}
The locked yield $L$ is set at inception $\IRR$, so $y-L$ is the 6M
$\IRR$ performance $R_{t+6M} - R_t$.
\begin{center}\small
\begin{tabular}{lcc}
\toprule
\textbf{6M} & $y - L$ & $(y-L)^2$\\
\midrule
Average & $-0.090\%$ & $0.004\%$\\
StDev & $0.663\%$ & $0.006\%$\\
5th pct & $-1.44\%$ & $3.92\!\times\!10^{-8}$\\
95th pct & $0.93\%$ & $0.0206\%$\\
\bottomrule
\end{tabular}
\end{center}
With $M\sim 150$ from \cref{prop:overhedge}: average overhedge error
$\sim 0.32\%$ of notional, 95th-percentile bound $\sim 1.55\%$
(materialised in the 2011 US-downgrade episode).

\subsection{Upper bound for negative gamma --- \cref{eq:gamma-bound}}
\begin{center}\small
\begin{tabular}{lccc}
\toprule
\textbf{Upper bound on $y - L$} & \textbf{2Y} & \textbf{5Y} & \textbf{10Y}\\
\midrule
Average & 22.47\% & 21.31\% & 11.44\%\\
StDev & 15.20\% & 0.12\% & 0.10\%\\
5th pct & 11.15\% & 21.11\% & 11.28\%\\
95th pct & 25.57\% & 21.51\% & 11.61\%\\
\bottomrule
\end{tabular}
\end{center}
A move of $> 11\%$ in the 10Y yield within a typical \Tlock\ tenor
ATM has never been recorded.

\subsection{Roll P\&L on the 10Y (trader's view, 3\% locked yield)}
\begin{center}\small
\begin{tabular}{lcccc}
\toprule
\textbf{10Y Roll P\&L} & \textbf{1M} & \textbf{3M} & \textbf{6M} & \textbf{12M}\\
\midrule
Average & $-0.1085\%$ & $-0.1069\%$ & $-0.1045\%$ & $-0.0996\%$\\
StDev   & $0.5028\%$  & $0.4953\%$  & $0.4841\%$  & $0.4615\%$\\
5th pct & $0.4749\%$  & $0.4677\%$  & $0.4568\%$  & $0.4352\%$\\
95th pct & $-1.0395\%$ & $-1.0238\%$ & $-1.0005\%$ & $-0.9530\%$\\
\bottomrule
\end{tabular}
\end{center}

\subsection{Roll P\&L on the 5Y (trader's view, 3\%)}
\begin{center}\small
\begin{tabular}{lcccc}
\toprule
\textbf{5Y Roll P\&L} & \textbf{1M} & \textbf{3M} & \textbf{6M} & \textbf{12M}\\
\midrule
Average & $-0.05\%$ & $-0.04\%$ & $-0.04\%$ & $-0.04\%$\\
StDev   & $0.22\%$  & $0.22\%$  & $0.21\%$  & $0.18\%$\\
5th pct & $0.24\%$  & $0.23\%$  & $0.22\%$  & $0.20\%$\\
95th pct& $-0.43\%$ & $-0.42\%$ & $-0.40\%$ & $-0.35\%$\\
\bottomrule
\end{tabular}
\end{center}

\subsection{Roll P\&L on the 2Y (trader's view, 3\%)}
\begin{center}\small
\begin{tabular}{lcccc}
\toprule
\textbf{2Y Roll P\&L} & \textbf{1M} & \textbf{3M} & \textbf{6M} & \textbf{12M}\\
\midrule
Average & $-0.07\%$ & $-0.07\%$ & $-0.06\%$ & $-0.05\%$\\
StDev   & $0.51\%$  & $0.51\%$  & $0.50\%$  & $0.50\%$\\
5th pct & $0.32\%$  & $0.31\%$  & $0.31\%$  & $0.30\%$\\
95th pct& $-0.52\%$ & $-0.51\%$ & $-0.49\%$ & $-0.45\%$\\
\bottomrule
\end{tabular}
\end{center}


\section{Valuation Adjustments}\label{sec:xva}

A \Tlock\ position carries the standard XVA components:

\begin{description}[leftmargin=1.6em,itemsep=2pt]
\item[CVA / DVA.] Counterparty and own-default risk. Computation is
run-of-the-mill, but a short \Tlock\ is subject to \emph{systemic
wrong-way counterparty risk}: in distressed corporate-market scenarios,
a flight-to-quality lowers Treasury yields exactly when the corporate
client is most exposed, raising the bank's exposure to that client.
Pucci argues this may be (partially) offset by a specular increase in
DVA --- it is hard to envision a regime where the corporate sector
suffers and banks thrive.

\item[FVA.] The natural hedge is a repo trade. Funding cost arises
either from the cash leg the bank posts to borrow the security, or from
the collateralisation mechanism if the repo is CSA-assisted or cleared.

\item[KVA.] Capital costs are minimal --- one of the reasons \Tlock s
are profitable for the dealer bank, though by the same token mark-ups
must remain low.
\end{description}

\textbf{Proxy underestimation.} Because the forward proxy is an
overhedge, it \emph{underestimates} the exposure of the short \Tlock,
hence underestimates the "contra" contributions (CVA, FVA, KVA).
The short tenor mitigates the consequences, since XVAs for ATM short
trades are normally a fraction of the carried exposure.


\section{Conclusions (Pucci 2019)}

\begin{enumerate}[leftmargin=1.6em,itemsep=2pt]
\item A long Standard \Tlock\ has negative convexity locally near $L$.
\item A trader can safely book a short \Tlock\ as a long proxy forward
      contract on the same security with the same expiry. The forward
      is a natural overhedge, generating a \emph{conservative} P\&L
      representation.
\item Delta hedging the position generates non-negative \emph{gamma
      profits}.
\item A Then-Current \Tlock\ can be booked as a proxy forward on the
      current on-the-run, with revision at each roll date generating
      statistically negligible P\&L.
\end{enumerate}


\section{Validation and Numerical Examples}\label{sec:validation}

This section is the bridge to a pricing-library implementation.

\subsection{Pucci's worked example --- 24 January 2019}\label{sec:worked}
\textbf{Use this verbatim as a regression test.}

\begin{center}\small
\renewcommand{\arraystretch}{1.18}
\begin{tabular}{p{5.8cm} l}
\toprule
\textbf{Parameter} & \textbf{Value} \\
\midrule
Trade date $t$ & 24 January 2019 \\
Underlying & US Treasury, ISIN \texttt{US9128285M81} \\
Coupon rate $c$ & 3.125\% \\
Maturity & 25 November 2028 \\
Expiry $t_e$ & 24 April 2019 (3-month tenor) \\
Locked yield $L$ & 2.717\% (ATM) \\
Spot $\IRR_t$ & 2.717\% \\
Spot dirty price $P^{\mathrm{mkt}}_t$ & 104.1055\% \\
Spot clean price & 103.4922\% \\
3M repo rate $\rrepo$ & 2.46\% \\
Forward $\IRR$ ($\IRR_t$ at $t_e$) & 2.53\% \\
Forward dirty price & 104.74\% \\
Forward clean price ($\Bf$) & 103.36\% \\
\bottomrule
\end{tabular}
\end{center}

\textbf{Validation checks.} A pricing library implementing the formulas
of this document must reproduce, on this dataset:
\begin{enumerate}[leftmargin=1.6em,itemsep=2pt]
\item The bond \emph{forward clean price} 103.36\% via
      \cref{eq:bond-forward} (Part I) and \cref{eq:fwd-repo} (Part II).
      \textbf{The two must agree} up to the simple-vs-continuous
      compounding convention (Part I uses simple, Part II continuous);
      switch via $\Drepo = e^{-\rrepo(\Tdel-T_0)}$ vs
      $\Drepo = 1/(1+\rrepo(\Tdel-T_0))$.
\item The forward yield $y_f = 2.53\%$ as the inverse of
      \cref{eq:forward-yield}: $P_{\Tdel}(y_f) = \Bf$.
\item The \Tlock\ NPV via all three Part I conventions (\cref{eq:npv-price},
      \cref{eq:npv-yield}, dirty-price): at $L = \IRR_t = 2.717\%$, the
      \emph{spot} NPV is exactly zero (ATM); the forward NPV is non-zero
      because of carry.
\item The overhedge plot of Figure 1 in Pucci's paper:
      payoff $g(y)$ via \cref{eq:payoff-pucci},
      forward approx via \cref{eq:K-approx},
      overhedge error via $R_1(y)$ of \cref{prop:overhedge}. The error
      curve is convex with minimum (zero) at $y=L$ and small positive
      values on either side.
\end{enumerate}

\subsection{Library-ready parameter block (pseudocode)}

\begin{verbatim}
# --- Trade & market state ---
trade_date     = 2019-01-24
expiry         = 2019-04-24
isin           = "US9128285M81"
coupon         = 0.03125              # annual
maturity       = 2028-11-25
notional       = 1.0                  # 100% of face
locked_yield   = 0.02717

# --- Calibration targets ---
spot_dirty     = 1.041055
spot_clean     = 1.034922
repo_3M        = 0.0246
fwd_irr        = 0.02530
fwd_dirty      = 1.0474
fwd_clean      = 1.0336

# --- Coupon dates in (trade_date, expiry] ---
#   (zero or one for a 3M T-Lock on a semi-annual UST; check)
#   For 11/25 + 5/25 schedule, no coupon falls in (1/24, 4/24]
\end{verbatim}

\subsection{Cross-formula consistency}

\textbf{Bond forward, simple vs continuous.}
Part I writes \cref{eq:bond-forward} with simple-compounded discount
factors $\Drepo(s,T) = 1/(1+\rrepo(T-s))$. Part II writes
\cref{eq:fwd-repo} with continuous-compounded discount.
For the 3M example:
\begin{itemize}[leftmargin=1.4em,itemsep=1pt]
\item simple: $(1 + 0.0246\cdot0.25) = 1.00615$
\item continuous: $e^{0.0246\cdot0.25} = 1.006159$
\end{itemize}
identical to 5 decimal places.

\textbf{$\PVone(y_f)$ vs spot $\PVone$.}
The forward-yield $\PVone$ in \cref{eq:pv01-fwd} uses the bond's
price-yield function \emph{at expiry} (settlement = $\Tdel$). The spot
$\PVone$ uses settlement $T_0$. The two differ by the carry between
$T_0$ and $\Tdel$. Library implementation MUST distinguish them
(\cref{rem:pv01-vs-fwd}).

\textbf{Two-curve discounting in Part I.}
$\Bf$ uses $\Drepo$; the final cash payoff uses $\Dtl$ (bond zero-yield
curve). These are different curves. In Pucci (Part II), the formal
arbitrage-free formula uses a single risk-free $D_{tt_e}$, justified by
the short tenor + US-government issuer. The practitioner two-curve
approach is the more general (and more careful) representation.

\subsection{Risk: what to compute and how to interpret}\label{sec:risk-impl}

\begin{description}[leftmargin=1.6em,itemsep=2pt]
\item[Delta vs forward price.] \cref{eq:delta-P}; negative for a long
      \Tlock\ at ATM; hedge by shorting the bond via reverse repo.
\item[Gamma vs forward price.] \cref{eq:gamma-local}; negative for
      $|y-L|$ small, bounded explicitly by \cref{eq:gamma-bound}.
\item[Repo bucket vega.] \cref{eq:fwd-repo} depends on $\rrepo$ at
      $t_e$ and at each coupon date $t_i\in(t,t_e]$. Bump each tenor
      independently for a repo-curve risk decomposition.
\item[PV01 / NDELTA of the underlier.] Compute and report
      as approximations (the anonymous paper's caveat:
      \cref{sec:risks-practitioner}). They are \emph{not} the
      sensitivity of the \Tlock\ itself --- those are
      \cref{eq:delta-P,eq:gamma-local}.
\end{description}

\subsection{Smoke tests}

\begin{enumerate}[leftmargin=1.6em,itemsep=2pt]
\item \textbf{ATM zero NPV.}
      At $\ytl = \IRR_t$ on the trade date the spot \Tlock\ NPV is zero
      (after netting carry). Test all three conventions.
\item \textbf{Parallel shift symmetry.}
      Shift the bond yield by $\pm\Delta y$; payoffs are linear at first
      order (sign by sign), overhedge error symmetric and quadratic.
\item \textbf{Sign convention.} $a=+1$ long, $a=-1$ short; long pays
      when $\IRR > L$. A long \Tlock\ + a long bond forward should be
      approximately delta-neutral.
\item \textbf{No-arbitrage with the forward proxy.}
      $\text{NPV}^{\Tlock}_{\text{long}} +
       \text{NPV}^{\text{fwd, short, K=Pte(L)}}$
      $= R_1(y)$ remainder, bounded by \cref{eq:rem-bound}.
\end{enumerate}


\appendix

\section{Simple Math: Chain Rule and Inverse Function}\label{sec:simple-math}

For differentiable $f,g$:
\begin{align}
(f\circ g)'  &= f' \cdot g', \label{eq:chain1}\\
(f\circ g)'' &= (f''\circ g)\cdot(g')^2 + (f'\circ g)\cdot g''.
\label{eq:chain2}
\end{align}
For an invertible differentiable $f$:
\begin{align}
(f^{-1})' &= \frac{1}{f'}, \label{eq:inv1}\\
(f^{-1})'' &= -\frac{f''}{(f')^3}. \label{eq:inv2}
\end{align}
Applied to $y = P^{-1}(P)$ (treating $P$ as the independent variable
when computing greeks w.r.t.\ the bond price):
\[
D_P[y] = 1/D_y[P], \qquad
\dDP{2}[y] = -\dDy{2}[P]/(D_y[P])^3.
\]
These are the substitutions used in \cref{sec:gamma}.


\section{Key Derivatives of the Bond Price}\label{sec:key-derivs}

Decomposing $P^{(c)}_t(y) = P^{(0)}_t(y) + c\,A_t(y)$ with
\[
P^{(0)}_t(y) = e^{-(t_n-t)y}, \qquad
A_t(y) = \sum_{t_i > t} \alpha_i\, e^{-(t_i-t)y},
\]
we have
\begin{align}
D_y\bigl[P^{(c)}_t\bigr] &= -(t_n-t)e^{-(t_n-t)y} + c\, D_y[A_t]
                          = D_y\bigl[P^{(0)}_t\bigr] + c\,D_y[A_t],
\label{eq:DyPdec}\\
\dDy{k}\bigl[P^{(c)}\bigr] &= (t-t_n)^k e^{-(t_n-t)y}
        + c\sum_{t_i>t}\alpha_i(t-t_i)^k e^{-(t_i-t)y}, \label{eq:DkP}\\
\dDy{k}[A_t] &= \sum_{t_i > t}\alpha_i(t-t_i)^k e^{-(t_i-t)y},
\label{eq:DkA}\\
\dDy{k}\bigl[P^{(c)}_t\bigr] &= \dDy{k}\bigl[P^{(0)}_t\bigr]
                                 + c\,\dDy{k}[A_t]. \label{eq:DkPdec}
\end{align}


\section{More Math: The Convexity-Inequality Proof}\label{sec:more-math}

Throughout, set
\[
q_i \equiv (t - t_i), \qquad Q_i \equiv \alpha_i\,e^{-(t_i-t)y}.
\]

\begin{lemma}\label{lem:positivity}
$\dDy{2}\bigl[P^{(0)}_t\bigr]\,\dDy{2}[A_t]
 + \dDy{2}\bigl[P^{(0)}_t\bigr]\,\dDy{2}[A_t]
 - 2\,\dDy{2}\bigl[P^{(0)}_t\bigr]\,\dDy{2}[A_t] \;\ge\; 0$.
\end{lemma}

\begin{proof}
After substitution this becomes
\[
q_n Q_n \sum_{t_i>t} q_i^3 Q_i
\;+\; q_n^3 Q_n \sum_{t_i>t} q_i Q_i
\;-\; 2 q_n^2 Q_n \sum_{t_i>t} q_i Q_i
\;=\;
\sum_{t_i>t} Q_i Q_n\, q_i q_n\,[q_i - q_n]^2
\;\ge\; 0,
\]
clearly non-negative.
\end{proof}

\begin{lemma}\label{lem:A-negativity}
$\bigl(\dDy{2}[A_t]\bigr)^2 - \dDy{2}[A_t]\cdot\dDy{2}[A_t] < 0$.
\end{lemma}

\begin{proof}
We must show
\[
\Bigl[\sum_{t_i>t} q_i Q_i\Bigr]
\Bigl[\sum_{t_i>t} q_i^3 Q_i\Bigr]
- \Bigl[\sum_{t_i>t} q_i^2 Q_i\Bigr]^2 \;>\; 0,
\]
since the LHS equals
$\sum_{t_i,t_j > t} Q_i Q_j\,[q_i - q_j]^2 > 0$
by Cauchy-Schwarz. Concretely:
\begin{align*}
\sum_i q_i Q_i \sum_j q_j^3 Q_j - \Bigl[\sum_i q_i^2 Q_i\Bigr]^2
&= \sum_{i \ne j} q_i q_j^3 Q_i Q_j - \sum_{i\ne j} q_i^2 q_j^2 Q_i Q_j \\
&= \sum_{i<j}\bigl(q_i q_j^3 + q_i^3 q_j - 2 q_i^2 q_j^2\bigr) Q_i Q_j\\
&= \sum_{i<j} q_i q_j\, (q_i - q_j)^2\, Q_i Q_j \;>\; 0.
\qedhere
\end{align*}
\end{proof}

\begin{proof}[Proof of \cref{lem:bracket}]
We must show $(\dDy{2}[P])^2 - D_y[P]\cdot\dDy{3}[P] < 0$.
Using the decomposition \cref{eq:DkPdec}, and expanding:
\begin{align}
(\dDy{2}[P])^2
  &= \bigl(\dDy{2}[P^{(0)}]\bigr)^2
   + 2c\,\dDy{2}[P^{(0)}]\,\dDy{2}[A]
   + c^2 \bigl(\dDy{2}[A]\bigr)^2, \notag \\
D_y[P]\,\dDy{3}[P]
  &= D_y[P^{(0)}]\,\dDy{3}[P^{(0)}]
   + c\bigl[D_y[P^{(0)}]\,\dDy{3}[A] + D_y[A]\,\dDy{3}[P^{(0)}]\bigr]
   + c^2\, D_y[A]\,\dDy{3}[A]. \notag
\end{align}
The leading $\bigl(\dDy{2}[P^{(0)}]\bigr)^2$ and
$-D_y[P^{(0)}]\,\dDy{3}[P^{(0)}]$ cancel (both are computed directly
from the zero-coupon factor and one verifies that
$\bigl(\dDy{2}[P^{(0)}]\bigr)^2 = D_y[P^{(0)}]\,\dDy{3}[P^{(0)}]$, both
equal to $(t_n-t)^4\, e^{-2(t_n-t)y}$). Subtracting and grouping:
\begin{equation}
(\dDy{2}[P])^2 - D_y[P]\,\dDy{3}[P]
\;=\; c\,\mathcal{L} \;+\; c^2\,\mathcal{Q},
\label{eq:bracket-proof}
\end{equation}
with
\begin{align*}
\mathcal{L} &= 2\,\dDy{2}[P^{(0)}]\,\dDy{2}[A]
              - D_y[P^{(0)}]\,\dDy{3}[A]
              - D_y[A]\,\dDy{3}[P^{(0)}], \\
\mathcal{Q} &= \bigl(\dDy{2}[A]\bigr)^2 - D_y[A]\,\dDy{3}[A].
\end{align*}
By \cref{lem:positivity} $\mathcal{L} \le 0$ and by \cref{lem:A-negativity}
$\mathcal{Q} < 0$. With $c > 0$, the whole expression is negative.
\end{proof}


\section{The Forward Price and the Repo Strategy}\label{sec:fwd-repo-deriv}

A \emph{coupon security} pays a discrete stream of cash flows $c_i$ at
times $T_i$. For a generic haircut $h^{\mathrm{cut}}\in[0,1]$, the
$T$-forward price under continuous compounding is
\begin{empheq}[box=\fbox]{align}
\FwdP_t(T) &= X_t\Bigl[(1-h^{\mathrm{cut}})e^{\rrepo(T-t)}
                       + h^{\mathrm{cut}}\,e^{\rfun(T-t)}\Bigr] \notag\\
           &\quad - \sum_{t < T_i \le T} c_i\,
                    e^{(h^{\mathrm{cut}}\rfun
                       + (1-h^{\mathrm{cut}})\rrepo)(T-T_i)}.
\label{eq:fwd-haircut}
\end{empheq}
Rates are continuously compounded and flat for ease of illustration.

\subsection{The repo-cash strategy (zero haircut)}

In a repo trade with intermediate coupons, the trader:
\begin{enumerate}[leftmargin=1.6em,itemsep=1pt]
\item At $t$, borrows $X_t$ from the repo counterparty at rate
      $\rrepo$ until $T$; buys the bond and transfers it to the
      counterparty as collateral.
\item On each coupon date $T_i$, receives coupon $c_i$ and decreases
      the outstanding by $c_i$.
\item At $T$, receives the bond back, sells it in the market, and
      returns all outstanding cash.
\end{enumerate}
Let $A_u$ = repo cash account at time $u$ (negative = owed). At start:
$A_t = -X_t$. At $T_1$:
\[
A_{T_1} = c_1 + A_t\, e^{\rrepo(T_1-t)}.
\]
Recursively for $i = 2,\dots,q$:
\[
A_{T_i} = c_i + A_{T_{i-1}}\,e^{\rrepo(T_i - T_{i-1})}.
\]
Between $T_q$ and $T$ it accrues to
\[
A_T = A_{T_q}\, e^{\rrepo(T-T_q)}.
\]

\begin{lemma}\label{lem:repo-account}
For all $u \le T$,
\begin{equation}
A_u \;=\; \sum_{t < T_i \le u} c_i\,e^{\rrepo(u - T_i)}
        \;-\; X_t\,e^{\rrepo(u - t)}.
\label{eq:Au}
\end{equation}
\end{lemma}

\begin{proof}
Induction. Base: $A_t = -X_t$. Inductive step from $T_{i-1}$ to $T_i$:
\begin{align*}
A_{T_i}
&= c_i + A_{T_{i-1}}\,e^{\rrepo(T_i - T_{i-1})} \\
&= c_i + \Bigl[\sum_{t<T_j\le T_{i-1}} c_j\,e^{\rrepo(T_{i-1}-T_j)}
              - X_t\,e^{\rrepo(T_{i-1}-t)}\Bigr]e^{\rrepo(T_i-T_{i-1})}\\
&= c_i + \sum_{t<T_j\le T_{i-1}} c_j\,e^{\rrepo(T_i - T_j)}
       - X_t\,e^{\rrepo(T_i - t)} \\
&= \sum_{t<T_j\le T_i} c_j\,e^{\rrepo(T_i - T_j)}
   - X_t\,e^{\rrepo(T_i - t)}.
\end{align*}
Free accrual between $T_q$ and $u\le T$ preserves the form. \qed
\end{proof}

\textbf{Forward-price identity.} The strategy is zero-cost at $t$ and
yields at $T$:
\[
v_T = X_T + A_T.
\]
By definition $A_T$ is the $T$-forward price:
\begin{empheq}[box=\fbox]{equation}
\FwdP_t(T) \;=\; X_t\,e^{\rrepo(T-t)}
                  - \sum_{t<T_i\le T} c_i\,e^{\rrepo(T - T_i)}.
\label{eq:fwd-zero-haircut}
\end{empheq}
This is \cref{eq:fwd-haircut} at $h^{\mathrm{cut}} = 0$.

\subsection{Unsecured-funding limit and the blend}

If the repo trade is unavailable, all funding is unsecured at rate
$\rfun$:
\begin{equation}
\FwdP_t(T) \;=\; X_t\,e^{\rfun(T-t)}
                  - \sum_{t<T_i\le T} c_i\,e^{\rfun(T - T_i)}.
\label{eq:fwd-full-funding}
\end{equation}
This is \cref{eq:fwd-haircut} at $h^{\mathrm{cut}} = 1$.

A generic haircut $0 \le h^{\mathrm{cut}} \le 1$ produces the blend
\cref{eq:fwd-haircut}, derivable from a strategy that distributes
coupon payments between repo redemption and unsecured-funding redemption
according to their relative weights.

For a Treasury we assume $h^{\mathrm{cut}} = 0$ throughout
the body of this study.


\section{References}\label{sec:bib}

\begin{enumerate}[leftmargin=2.2em,itemsep=1pt]
\item \textit{Treasury Lock Model.} Anonymous practitioner note, no
      author / no date.
      (The paper summarised in Part I.)
\item Pucci, M. \textit{Hedging the Treasury Lock.} Banca IMI SpA,
      Milan, August 19, 2019. SSRN no.\ 3386521.
\item Pucci, M. \textit{Constant Maturity Treasury Convexity Correction.}
      International Journal of Theoretical and Applied Finance, 17(8),
      2014. SSRN 3387961.
\item Adams, J. \& Smith, D.J. \textit{Pre-Issuance Hedging of
      Fixed-Rate Debt.} J.\ Applied Corporate Finance 23(4), 2011.
\item Hunt, P.J. \& Kennedy, J.E. \textit{Financial Derivatives in
      Theory and Practice.} Wiley, 2004.
\item Hull, J.C. \textit{Options, Futures and Other Derivatives.}
      Prentice Hall, 9th ed., 2015.
\item Taleb, N. \textit{Dynamic Hedging.} Wiley, 1996.
\item Carr, P. \& Madan, D. \textit{Towards a Theory of Volatility
      Trading.} 2002.
\item Ramirez, J. \textit{Accounting for Derivatives: Advanced Hedging
      under IFRS.} Wiley, 2015. \texttt{[6, 7.9]}.
\item PwC. \textit{In depth: Achieving hedge accounting in practice
      under IFRS 9.} 2017.
\item Gregory, J. \textit{Counterparty Credit Risk.} Wiley, 2010.
\item Morini, M. \& Prampolini, A. \textit{Risky Funding: A Unified
      Framework for Counterparty and Liquidity Charges.}
      Landmarks in XVA, 2016.
\item Andersen, L., Duffie, D. \& Song, Y. \textit{Funding Value
      Adjustments.} Journal of Finance 74(1):145--192, 2019.
\item ISDA. \textit{Eight Treasury Rate Lock Agreements.} U.S.\ SEC,
      2005.
\item ICMA Group. \textit{Frequently Asked Questions on Repo.} 2015.
\end{enumerate}

\end{document}
